%! Characteristics of Linear Functions — Unit 2, Lesson 2.0 — Exit Ticket (Answer Key)
\documentclass[10pt]{article}
\usepackage{algebra2-article}
\usepackage{algebra2-key}   % pulls in -boxes; do NOT also load -boxes
\newcommand{\TallMath}[1]{$\displaystyle #1\rule[-1.4em]{0pt}{3.2em}$}
\newcommand{\lgrid}[4]{%
  \draw[step=1, linegray!40, very thin] (#1,#3) grid (#2,#4);
  \draw[->, semithick, charcoal] (#1,0)--(#2,0) node[below right, font=\scriptsize]{$x$};
  \draw[->, semithick, charcoal] (0,#3)--(0,#4) node[left, font=\scriptsize]{$y$};}

\begin{document}

\pageheader{Unit 2, Lesson 2.0}{Exit Ticket --- Answer Key}

\vspace{0.10in}

No notes. The graph below is the line $f(x) = -2x + 4$. Use it for all three questions.
\begin{center}
\begin{tikzpicture}[scale=0.52]
  \lgrid{-2}{5}{-3}{6}
  \draw[very thick, forest] (-1,6)--(3.5,-3);
  \fill[charcoal] (0,4) circle (3pt);
  \fill[charcoal] (2,0) circle (3pt);
\end{tikzpicture}
\end{center}

\begin{enumerate}[leftmargin=1.8em, itemsep=12pt]
  \item \textbf{Read the features.} Give the $y$-intercept, the $x$-intercept (zero), and the slope.
        \ansline{$y$-intercept $(0,4)$; \; $x$-intercept / zero $(2,0)$; \; slope $m=-2$ (down $2$ for every right $1$).}
  \item \textbf{Classify.} Is $f$ increasing or decreasing? State the domain and the range.
        \ansline{Decreasing (slope $<0$). Domain: all real numbers $(-\infty,\infty)$; range: all real numbers $(-\infty,\infty)$.}
  \item \textbf{Justify.} Name the interval where $f(x) > 0$ (the function is positive). Then explain,
        using the \emph{zero}, why the sign of $f$ switches there.
        \ansline{$f(x)>0$ for $x<2$, i.e.\ on $(-\infty,2)$. At the zero $x=2$, $f(2)=0$; left of it the graph is above the $x$-axis (positive), right of it below (negative). A line crosses the $x$-axis only once, so the sign changes exactly at the zero.}
\end{enumerate}

\end{document}
